\section*{Лекция 04 - 28 октомври}

\section{Подредици и точки на сгъстяване. Теорема на Болцано-Вайерщрас (принцип за компактност)}

Понятията „точка на сгъстяване“ и „подредица“ са тясно свързани, както ще се убедим след малко. Важно е да си дадете точна сметка за разликата между „граница“ и „точка на сгъстяване“.

\begin{definition}{Точка на сгъстяване на редица}{seq-cluster-point}
Нека $\{a_n\}_{n=1}^{\infty} \subset \mathbb{R}$ и $a \in \mathbb{R}$. Казваме, че $a$ е \textbf{точка на сгъстяване} на $\{a_n\}_{n=1}^{\infty}$, ако във всяка околност на $a$ има безброй много членове на редицата. Формално, за всяко $U$ -- околност на точката $a$, множеството
\[
\{n \in \mathbb{N} : a_n \in U\}
\]
е безкрайно.
\end{definition}

\noindent\textbf{Забележка.} Забележете, че броим индексите! Помислете за точките на сгъстяване на $0, 1, 0, 1, \dots$.

\begin{theorem}{Единственост на точката на сгъстяване при сходяща редица}{prop-convergent-unique-cluster}
Ако $a_n \xrightarrow[n\to\infty]{} a$, то $a$ е точка на сгъстяване на редицата $\{a_n\}_{n=1}^{\infty}$, при това единствена.
\end{theorem}

\begin{proof}
Границата е точка на сгъстяване, защото кофинитните множества винаги са безкрайни. Допускаме, че съществува друга точка на сгъстяване $b \in \mathbb{R}$, $a \ne b$. Ако $U$ е околност на $a$ и $V$ е околност на $b$, такива че $U \cap V = \emptyset$, то можем да приложим дефинициите за сходяща редица и да достигнем до противоречие с дефиницията за точка на сгъстяване:
\[
\underbrace{\{n \in \mathbb{N} : a_n \in U\}}_{\text{кофинитно}} \implies \underbrace{\{n \in \mathbb{N} : a_n \notin U\}}_{\text{крайно}} \supset \underbrace{\{n \in \mathbb{N} : a_n \in V\}}_{\text{безкрайно}}.
\]
Очевидно безкрайно множество няма как да се съдържа в крайно множество. Включването се дължи на факта, че $U \cap V = \emptyset \implies \mathbb{R} \setminus U \supset V$.
\end{proof}

\paragraph{Пример 1.3.}
Като пример можем да разгледаме $a_n \xrightarrow[n\to\infty]{} a$, $b_n \xrightarrow[n\to\infty]{} b$ ($a \ne b$) и редицата $a_1, b_1, a_2, b_2, \dots, a_n, b_n, \dots$. Лесно се съобразява, че именно границите $a$ и $b$ на горните редици са точки на сгъстяване на новата редица. Можем да посочим и растяща редица, която няма точки на сгъстяване: $1, 2, 3, \dots, n, \dots$.

\paragraph{Пример 1.4.}
Знаем, че множеството от рационалните числа $\mathbb{Q}$ е изброимо и следователно може да бъде подредено в редица, нека например
\[
\mathbb{Q} = \{q_1, q_2, \dots, q_n, \dots\}.
\]
Съобразете, че всяко реално число е точка на сгъстяване на горната редица.

\begin{definition}{Подредица}{subsequence}
Ако от редица $\{a_n\}_{n=1}^{\infty}$ задраскаме част от членовете така, че да останат безброй много членове и запазим реда на останалите, получаваме \textbf{подредица} на първоначалната. Еквивалентно, всяка строго растяща (безкрайна) редица от естествени числа $n_1 < n_2 < n_3 < \dots < n_k < \dots$ задава подредица $\{a_{n_k}\}_{k=1}^{\infty}$ на $\{a_n\}_{n=1}^{\infty}$.
\end{definition}

Да си спомним дефиницията на редица като изображение $a \colon \mathbb{N} \to \mathbb{R}$. Сега казваме, че всяка нейна подредица е композиция на някакво строго растящо изображение $n \colon \mathbb{N} \to \mathbb{N}$ и на $a \colon \mathbb{N} \to \mathbb{R}$:
\[
k \mapsto n_k \mapsto a_{n_k}.
\]

\begin{theorem}{Сходимост на подредица}{prop-subseq-convergence}
Нека $\{a_n\}_{n=1}^{\infty}$ е сходяща редица от реални числа и $\{a_{n_k}\}_{k=1}^{\infty}$ е нейна подредица. Тогава $\{a_{n_k}\}_{k=1}^{\infty}$ е сходяща към границата на $\{a_n\}_{n=1}^{\infty}$.
\end{theorem}

\begin{proof}
Нека $a_n \xrightarrow[n\to\infty]{} a$ и $U$ е произволна околност на $a$. Тогава съществува $n_0 \in \mathbb{N}$ такова, че за всяко $n \ge n_0$ е в сила $a_n \in U$. Тъй като $n_1 < n_2 < \dots < n_k < \dots$ е строго растяща и безкрайна, съществува $k_0 \in \mathbb{N}$ такова, че $n_{k_0} \ge n_0$. Тогава за всяко $k \ge k_0$ имаме $n_k \ge n_{k_0} \ge n_0$ и следователно $a_{n_k} \in U$. Получихме, че $a_{n_k} \xrightarrow[k\to\infty]{} a$.
\end{proof}

\begin{theorem}{Връзка между подредици и точки на сгъстяване}{prop-subseq-cluster-equiv}
Нека $\{a_n\}_{n=1}^{\infty} \subset \mathbb{R}$ и $a \in \mathbb{R}$. Тогава $a$ е точка на сгъстяване на дадената редица точно тогава, когато съществува подредица $\{a_{n_k}\}_{k=1}^{\infty}$ на $\{a_n\}_{n=1}^{\infty}$, за която е изпълнено:
\[
a_{n_k} \xrightarrow[k\to\infty]{} a.
\]
\end{theorem}

\begin{proof}
$(\Leftarrow)$ Ако $\{a_{n_k}\}_{k=1}^{\infty}$ е подредица на $\{a_n\}_{n=1}^{\infty}$ и при това $a_{n_k} \xrightarrow[k\to\infty]{} a$, то е ясно, че $a$ е точка на сгъстяване на $\{a_n\}_{n=1}^{\infty}$. Наистина, ако $U$ е околност на $a$, то:
\[
\{n \in \mathbb{N} : a_n \in U\} \supset \{n_k = n(k) \in \mathbb{N} : k \ge k_0\},
\]
защото след определен индекс $k_0$ всички членове на подредицата попадат в избраната околност на $a$.

$(\Rightarrow)$ Обратно, нека $a$ е точка на сгъстяване за $\{a_n\}_{n=1}^{\infty}$. Интервалът $(a - 1, a + 1)$ е околност на $a$. Тогава съществува $n_1 \in \mathbb{N}$ такова, че $a_{n_1} \in (a - 1, a + 1)$. Разглеждаме $\left(a - \frac{1}{2}, a + \frac{1}{2}\right)$ -- околност на $a$. Понеже $\{n \in \mathbb{N} : a_n \in (a - 1, a + 1)\}$ е безкрайно (от дефиницията за точка на сгъстяване) и $\{1, 2, \dots, n_1\}$ е крайно, то съществува $n_2 > n_1$ с $a_{n_2} \in \left(a - \frac{1}{2}, a + \frac{1}{2}\right)$.

Продължаваме с аналогични разсъждения и избираме индекси $n_1 < n_2 < \dots < n_{k-1}$. На поредната стъпка имаме околност $\left(a - \frac{1}{k}, a + \frac{1}{k}\right)$ на $a$. Понеже множеството
\[
\left\{n \in \mathbb{N} : a_n \in \left(a - \frac{1}{k}, a + \frac{1}{k}\right)\right\}
\]
е безкрайно и $\{1, 2, \dots, n_{k-1}\}$ е крайно, то съществува $n_k > n_{k-1}$ такова, че $a_{n_k} \in \left(a - \frac{1}{k}, a + \frac{1}{k}\right)$.

По този начин построихме $n_1 < n_2 < n_3 < \dots < n_k < \dots$ (и следователно $\{a_{n_k}\}_{k=1}^{\infty}$ е подредица на $\{a_n\}_{n=1}^{\infty}$) така, че
\[
|a_{n_k} - a| < \frac{1}{k} \xrightarrow[k\to\infty]{} 0 \implies a_{n_k} \xrightarrow[k\to\infty]{} a.
\]
С това доказателството е завършено.
\end{proof}

\begin{theorem}{Теорема на Болцано--Вайерщрас (Принцип за компактност)}{bolzano-weierstrass}
Всяка ограничена редица има точка на сгъстяване. Еквивалентно, всяка ограничена редица има сходяща подредица.
\end{theorem}

\begin{proof}
Първо да съобразим, че ограничеността на $\{a_n\}_{n=1}^{\infty}$ влече съществуването на $b, c \in \mathbb{R}$ такива, че $a_n \in [b, c]$ за всяко $n \in \mathbb{N}$. Да положим $b_0 := b$, $c_0 := c$. Разглеждаме интервалите $\left[b_0, \frac{b_0 + c_0}{2}\right]$ и $\left[\frac{b_0 + c_0}{2}, c_0\right]$. Тогава поне едно от следните множества е безкрайно:
\[
\left\{n \in \mathbb{N} : a_n \in \left[b_0, \frac{b_0 + c_0}{2}\right]\right\}, \qquad \left\{n \in \mathbb{N} : a_n \in \left[\frac{b_0 + c_0}{2}, c_0\right]\right\}.
\]
Ако първото от тях е безкрайно, избираме $b_1 := b_0$, $c_1 := \frac{b_0 + c_0}{2}$. Ако първото множество е крайно, второто множество задължително е безкрайно и тогава полагаме $b_1 := \frac{b_0 + c_0}{2}$, $c_1 := c_0$. Съсредоточаваме се в интервала $[b_1, c_1]$, в който има безброй членове на редицата, и продължаваме по аналогичен начин.

Нека за някакво $k \in \mathbb{N} \cup \{0\}$ сме построили интервалите
\[
[b_0, c_0] \supset [b_1, c_1] \supset [b_2, c_2] \supset \dots \supset [b_k, c_k]
\]
такива, че множеството $\{n \in \mathbb{N} : a_n \in [b_k, c_k]\}$ е безкрайно и при това
\[
c_i - b_i = \frac{c - b}{2^i} \quad \forall i, \; 0 \le i \le k.
\]
Индукционната стъпка е естествена -- разглеждаме множествата:
\[
\left\{n \in \mathbb{N} : a_n \in \left[b_k, \frac{b_k + c_k}{2}\right]\right\}, \qquad \left\{n \in \mathbb{N} : a_n \in \left[\frac{b_k + c_k}{2}, c_k\right]\right\}.
\]
Ако първото от тези множества е безкрайно, полагаме $b_{k+1} := b_k$, $c_{k+1} := \frac{b_k + c_k}{2}$. Ако не, второто от двете множества е безкрайно (поради индукционното предположение) и тогава полагаме $b_{k+1} := \frac{b_k + c_k}{2}$, $c_{k+1} := c_k$. Сега имаме $[b_k, c_k] \supset [b_{k+1}, c_{k+1}]$, множеството $\{n \in \mathbb{N} : a_n \in [b_{k+1}, c_{k+1}]\}$ е безкрайно и при това:
\[
c_{k+1} - b_{k+1} = \frac{c_k - b_k}{2} = \frac{c - b}{2^{k+1}}.
\]
В тази конструкция избирахме втория интервал само ако в първия има краен брой членове на редицата. Направихме това за определеност: ако и в двата интервала има безброй членове на редицата, спокойно можем да изберем кой да е от тях и индукционното предположение пак ще бъде в сила за новата стъпка.

И тъй, построихме редица от вложени един в друг интервали:
\[
[b_0, c_0] \supset [b_1, c_1] \supset [b_2, c_2] \supset \dots \supset [b_k, c_k] \supset [b_{k+1}, c_{k+1}] \supset \dots
\]
Имаме
\[
\begin{aligned}
b_0 \le b_1 \le b_2 \le \dots \le b_k \le b_{k+1} \le \dots \\
c_0 \ge c_1 \ge c_2 \ge \dots \ge c_k \ge c_{k+1} \ge \dots
\end{aligned}
\]
Редиците $\{b_k\}_{k=1}^{\infty}$ и $\{c_k\}_{k=1}^{\infty}$ се схождат към едно и също число $a$. Наистина, тъй като $\{b_k\}_{k=1}^{\infty}$ е ограничена отгоре (от $c_0$) и монотонно растяща, тя е сходяща. Да означим границата ѝ с $a$: $b_k \xrightarrow[k\to\infty]{} a$. От друга страна:
\[
|c_k - a| \le |c_k - b_k| + |b_k - a| = \frac{c - b}{2^k} + |b_k - a| \xrightarrow[k\to\infty]{} 0 + 0 \implies c_k \xrightarrow[k\to\infty]{} a.
\]
Твърдим, че $a$ е точка на сгъстяване на $\{a_n\}_{n=1}^{\infty}$. Нека $\varepsilon > 0$ е произволно. Тъй като $(a - \varepsilon, a + \varepsilon)$ е околност на $a$, имаме:
\[
\left.
\begin{aligned}
b_k \xrightarrow[k\to\infty]{} a \implies b_k \in (a - \varepsilon, a + \varepsilon) \quad \forall k \ge k_1 \\
c_k \xrightarrow[k\to\infty]{} a \implies c_k \in (a - \varepsilon, a + \varepsilon) \quad \forall k \ge k_2
\end{aligned}
\right\} \implies \text{за } k_0 := \max\{k_1, k_2\}
\]
Имаме
\[
\begin{cases}
b_{k_0} \in (a - \varepsilon, a + \varepsilon) \\
c_{k_0} \in (a - \varepsilon, a + \varepsilon)
\end{cases}
\]
Сега $a - \varepsilon < b_{k_0} < c_{k_0} < a + \varepsilon$. Следователно според конструкцията на интервалите $[b_k, c_k]$, множеството $\{n \in \mathbb{N} : a_n \in [b_{k_0}, c_{k_0}]\}$ е безкрайно и остава да съобразим, че
\[
\{n \in \mathbb{N} : a_n \in [b_{k_0}, c_{k_0}]\} \subset \{n \in \mathbb{N} : a_n \in (a - \varepsilon, a + \varepsilon)\}.
\]
Това доказва, че $a$ е точка на сгъстяване на редицата $\{a_n\}_{n=1}^{\infty}$.
\end{proof}

\noindent\textbf{Важна забележка.} От горното доказателство директно се получава, че ако $\{a_n\}_{n=1}^{\infty}$ е редица от реални числа и $b, c \in \mathbb{R}$ са такива, че $\{n \in \mathbb{N} : a_n \in [b, c]\}$ е безкрайно, то $\{a_n\}_{n=1}^{\infty}$ има точка на сгъстяване, която принадлежи на интервала $[b, c]$.

\paragraph{Пример 1.9.}
Редицата
\[
1, \frac{1}{2}, 2, \frac{1}{3}, 3, \frac{1}{4}, 4, \dots
\]
има единствена точка на сгъстяване $0$, но не е ограничена и следователно не е сходяща. Да си припомним, че в Твърдение~\ref{thm:prop-convergent-unique-cluster} доказахме, че сходящите редици имат единствена точка на сгъстяване. Примерът показва, че обратното твърдение не е вярно в общия случай.

\begin{theorem}{Сходимост на ограничена редица с единствена точка на сгъстяване}{prop-bounded-unique-cluster}
Нека $\{a_n\}_{n=1}^{\infty}$ е ограничена редица от реални числа, която има единствена точка на сгъстяване. Тогава редицата е сходяща.
\end{theorem}

\begin{proof}
Имаме, че $b \le a_n \le c$ за всяко $n \in \mathbb{N}$, тъй като редицата е ограничена. Нека $\varepsilon > 0$ и разгледаме околност $(a - \varepsilon, a + \varepsilon)$ на $a$ -- единствената точка на сгъстяване на $\{a_n\}_{n=1}^{\infty}$. Ще докажем, че $a$ е граница на $\{a_n\}_{n=1}^{\infty}$.

Допускаме, че това не е вярно. Тогава съществува $\varepsilon > 0$ такова, че извън $(a - \varepsilon, a + \varepsilon)$ има безброй много членове на редицата, т.е. че $\{n \in \mathbb{N} : a_n \notin (a - \varepsilon, a + \varepsilon)\}$ е безкрайно. Следователно или $\{n \in \mathbb{N} : a_n \in [b, a - \varepsilon]\}$, или $\{n \in \mathbb{N} : a_n \in [a + \varepsilon, c]\}$ е безкрайно. (Ако например $a + \varepsilon > c$, приемаме, че $[a + \varepsilon, c] = \emptyset$.)

Използваме важната забележка, за да получим:
\begin{itemize}
    \item Ако първото от двете множества е безкрайно, то съществува точка на сгъстяване $d$ на $\{a_n\}_{n=1}^{\infty}$ такава, че $d \in [b, a - \varepsilon] \implies d \ne a$.
    \item Ако второто от двете множества е безкрайно, то съществува точка на сгъстяване $d'$ на $\{a_n\}_{n=1}^{\infty}$ такава, че $d' \in [a + \varepsilon, c] \implies d' \ne a$.
\end{itemize}
И в двата случая получаваме противоречие с единствеността на точката на сгъстяване $a$.
\end{proof}

\section{Необходимо и достатъчно условие на Коши за сходимост на числови редици}

Да си спомним, че една редица $\{a_n\}_{n=1}^{\infty}$ е сходяща, ако съществува реално число $a$ такова, че $a_n \xrightarrow[n\to\infty]{} a$. В тази дефиниция освен редицата участва и границата, за която предполагаме, че е неизвестна. Добре е да разполагаме с условие за сходимост, в което участват само членовете на редицата, но не и евентуалната граница.

\begin{definition}{Фундаментална редица (Редица на Коши)}{fundamental-sequence}
Казваме, че редицата $\{a_n\}_{n=1}^{\infty}$ е \textbf{фундаментална}, ако е изпълнено:
\[
\forall \varepsilon > 0 \; \exists n_0 \in \mathbb{N} \; \forall m \ge n_0 \; \forall n \ge n_0 \colon |a_m - a_n| < \varepsilon.
\]
\end{definition}

Необходимото и достатъчно условие на Коши за сходимост на числова редица е друго условие, което по същество е еквивалентно на принципа на непрекъснатост.

\begin{theorem}{Необходимо и достатъчно условие на Коши за сходимост на числова редица}{cauchy-criterion-seq}
Една редица от реални числа е сходяща точно тогава, когато е фундаментална.
\end{theorem}

\begin{proof}
Доказателството провеждаме в двете посоки:

$(\Rightarrow)$ \textbf{$\{a_n\}_{n=1}^{\infty}$ е сходяща $\implies \{a_n\}_{n=1}^{\infty}$ е фундаментална:}

Нека $a_n \xrightarrow[n\to\infty]{} a$ и $\varepsilon > 0$ е произволно. Тогава $\frac{\varepsilon}{2} > 0$ и следователно съществува $n_0 \in \mathbb{N}$ такова, че за всички $n \ge n_0$ е в сила $|a_n - a| < \frac{\varepsilon}{2}$. Да вземем кои да е две числа $m \ge n_0$ и $n \ge n_0$. Тогава:
\[
\left.
\begin{aligned}
|a_m - a| < \frac{\varepsilon}{2} \\
|a_n - a| < \frac{\varepsilon}{2}
\end{aligned}
\right\} \implies |a_m - a_n| = |a_m - a + a - a_n| \le |a_m - a| + |a_n - a| < \frac{\varepsilon}{2} + \frac{\varepsilon}{2} = \varepsilon.
\]

$(\Leftarrow)$ \textbf{$\{a_n\}_{n=1}^{\infty}$ е фундаментална $\implies \{a_n\}_{n=1}^{\infty}$ е сходяща:}

Преди всичко да покажем, че от фундаменталността на една редица следва нейната ограниченост. Наистина, да вземем $\varepsilon = 15 > 0$. Тогава съществува $n_0 \in \mathbb{N}$ такова, че за всички $n \ge n_0$ и за всички $m \ge n_0$ е в сила $|a_n - a_m| < 15$. В частност, при $m = n_0$, получаваме, че за всички $n \ge n_0$ е в сила $|a_n - a_{n_0}| < 15$. Следователно
\[
a_n \in (a_{n_0} - 15, a_{n_0} + 15) \quad \forall n \ge n_0.
\]
Избираме:
\[
\begin{cases}
b := \min \{a_1, a_2, a_3, \dots, a_{n_0-1}, a_{n_0} - 15\} \implies b \le a_n \quad \forall n \in \mathbb{N}, \\
c := \max \{a_1, a_2, a_3, \dots, a_{n_0-1}, a_{n_0} + 15\} \implies c \ge a_n \quad \forall n \in \mathbb{N}.
\end{cases}
\]
Следователно $\{a_n\}_{n=1}^{\infty}$ е ограничена.

Съгласно принципа за компактност (Теорема на Болцано--Вайерщрас) съществува подредица $\{a_{n_k}\}_{k=1}^{\infty}$ на $\{a_n\}_{n=1}^{\infty}$, която е сходяща. Нека $a_{n_k} \xrightarrow[k\to\infty]{} a$. Ще докажем, че цялата редица клони към $a$, т.е. $a_n \xrightarrow[n\to\infty]{} a$.

Избираме произволно $\varepsilon > 0$. Идеята е да използваме подходящ член на подредицата, за да направим числото
\[
|a_n - a| = |a_n - a_{n_k} + a_{n_k} - a| \le |a_n - a_{n_k}| + |a_{n_k} - a|
\]
по-малко от $\varepsilon$ за сметка на индекса $n$.

От фундаменталността на редицата имаме, че съществува $n_0 \in \mathbb{N}$ такова, че за всички $n \ge n_0$ и за всички $m \ge n_0$ е в сила $|a_n - a_m| < \frac{\varepsilon}{2}$. Тогава, ако $k$ е толкова голямо, че $n_k \ge n_0$, за всички $n \ge n_0$ ще получим $|a_n - a_{n_k}| < \frac{\varepsilon}{2}$, тоест ще направим първото събираемо малко. Тъй като $n_k \to \infty$, съществува $k_1 \in \mathbb{N}$ такова, че за всички $k \ge k_1$ е в сила $n_k \ge n_0$.

За да направим второто събираемо малко, използваме факта, че подредицата клони към $a$. От $a_{n_k} \xrightarrow[k\to\infty]{} a$ и $\frac{\varepsilon}{2} > 0$ следва, че съществува $k_2 \in \mathbb{N}$ такова, че за всички $k \ge k_2$ е в сила $|a_{n_k} - a| < \frac{\varepsilon}{2}$.

Фиксираме произволно $n \ge n_0$ ($n_0$ зависи само от $\varepsilon$). Избираме $k := \max\{k_1, k_2\}$. Тогава от $k \ge k_1$ имаме $|a_n - a_{n_k}| < \frac{\varepsilon}{2}$, а от $k \ge k_2$ имаме $|a_{n_k} - a| < \frac{\varepsilon}{2}$. Следователно:
\[
|a_n - a| \le |a_n - a_{n_k}| + |a_{n_k} - a| < \frac{\varepsilon}{2} + \frac{\varepsilon}{2} = \varepsilon
\]
за всички $n \ge n_0$. С това доказателството е завършено.
\end{proof}

\paragraph{Пример 2.3.}
Да разгледаме редицата с общ член $a_n = 1 + \frac{1}{2} + \dots + \frac{1}{n}$. Ясно е, че тази редица е строго растяща ($a_{n+1} = a_n + \frac{1}{n+1} > a_n$). При това:
\[
a_{2^{k+1}} - a_{2^k} = \frac{1}{2^k + 1} + \frac{1}{2^k + 2} + \dots + \frac{1}{2^k + 2^k} > 2^k \cdot \frac{1}{2^{k+1}} = \frac{1}{2}.
\]
Виждаме, че $\{a_n\}_{n=1}^{\infty}$ не е фундаментална (няма как да удовлетворява условието за фундаменталност при $\varepsilon \le \frac{1}{2}$). Следователно $\{a_n\}_{n=1}^{\infty}$ не може да бъде сходяща. Като използваме, че редицата е растяща, получаваме и че тя не е ограничена отгоре.

\chapter{Граници на функции}

Следващата ни цел е да въведем понятието „граница на функция, когато аргументът клони към $x_0$“ ($\lim_{x\to x_0} f(x)$). Интуицията би трябвало да бъде какво е поведението на втората координата на точка от графиката на функцията, когато първата координата се приближава към $x_0$.

\paragraph{Пример 3.1.}
\begin{itemize}
    \item Дефиниционната област на $f(x) = \frac{x^2 - 1}{x - 1}$ е $\operatorname{Dom}(f) = \mathbb{R} \setminus \{1\}$. Графиката на тази функция е права (графиката на $g(x) = x + 1$), от която е извадена точката с координати $(1, 2)$. Ако една точка е върху тази графика и първата ѝ координата се приближава към $1$, втората ѝ координата се приближава към $2$: $\lim_{x\to 1} f(x) = 2$.
    \item Да разгледаме функцията „цяла част на $x$“: $f(x) = \lfloor x \rfloor$. Тя е дефинирана върху цялата реална права. Поведението на втората координата на точка от нейната графика, когато първата координата се приближава към цяло число, е различно в зависимост от това дали първата координата се приближава отляво или отдясно. Съответно $\lim_{x\to x_0} f(x)$ не съществува, когато $x_0$ е цяло число.
    \item Функцията $f(x) = \sin \frac{1}{x}$ е с дефиниционна област $\operatorname{Dom}(f) = \mathbb{R} \setminus \{0\}$. Тя няма граница при $x \to 0$ по по-сложна причина: ако изберем колкото искаме малка околност на нулата, върху графиката има точки с първа координата в тази околност и втора координата кое да е реално число в интервала $[-1, 1]$.
\end{itemize}

\section{Точки на сгъстяване на множество}
Преди да въведем строго понятието „граница на функция, когато аргументът клони към $x_0$“, трябва да имаме представа кога аргументът на функция с дефиниционна област $D \subset \mathbb{R}$ може „да се приближава към $x_0$“.

\begin{definition}{Точка на сгъстяване на множество}{set-cluster-point}
Нека $D \subset \mathbb{R}$ и $x_0 \in \mathbb{R}$. Казваме, че $x_0$ е \textbf{точка на сгъстяване} на $D$, ако за всяка околност $U$ на $x_0$ множеството $U \cap D$ е безкрайно.
\end{definition}

\paragraph{Пример 3.3.}
\begin{itemize}
    \item Множеството от точките на сгъстяване на интервала $(3, +\infty)$ е интервалът $[3, +\infty)$.
    \item Множеството от точките на сгъстяване на $D = (-\infty, -1] \cup (0, 1) \cup \{5\} \cup \left\{17 + \frac{1}{n} : n \in \mathbb{N}\right\}$ е $(-\infty, -1] \cup [0, 1] \cup \{17\}$.
    \item Правете разлика между точка на сгъстяване на редица от реални числа и точка на сгъстяване на множество! Например точките на сгъстяване на редицата $0, 1, 0, 1, \dots$ са $0$ и $1$, а множеството от нейните стойности $\{0, 1\}$ няма точки на сгъстяване.
\end{itemize}

\begin{theorem}{Еквивалентни свойства за точка на сгъстяване на множество}{prop-cluster-point-set-equiv}
Нека $D \subset \mathbb{R}$, $x_0 \in \mathbb{R}$. Тогава следните твърдения са еквивалентни:
\begin{enumerate}
    \item $x_0$ е точка на сгъстяване на $D$.
    \item Във всяка околност $U$ на $x_0$ има точка от $D$, различна от $x_0$ ($(U \cap D) \setminus \{x_0\} \ne \emptyset$).
    \item Съществува редица $\{x_n\}_{n=1}^{\infty} \subset D$, $x_n \ne x_0$ за всяко $n \in \mathbb{N}$, такава че $x_n \xrightarrow[n\to\infty]{} x_0$.
    \item Съществува редица $\{x_n\}_{n=1}^{\infty} \subset D \setminus \{x_0\}$ такава, че $x_n \xrightarrow[n\to\infty]{} x_0$ и $|x_{n+1} - x_0| < |x_n - x_0|$ за всяко $n \in \mathbb{N}$.
\end{enumerate}
\end{theorem}

\begin{proof}
\textbf{(а) $\implies$ (б):} Тривиално, тъй като безкрайно множество не може да се състои единствено от $\{x_0\}$.

\textbf{(б) $\implies$ (в):} За произволно $n \in \mathbb{N}$ разглеждаме околност $\left(x_0 - \frac{1}{n}, x_0 + \frac{1}{n}\right)$ на $x_0$:
\[
\left[ D \cap \left(x_0 - \frac{1}{n}, x_0 + \frac{1}{n}\right) \right] \setminus \{x_0\} \ne \emptyset \implies \exists x_n \in D, \; x_n \ne x_0, \; |x_n - x_0| < \frac{1}{n}.
\]
Построихме $\{x_n\}_{n=1}^{\infty} \subset D \setminus \{x_0\}$ с $|x_n - x_0| < \frac{1}{n} \xrightarrow[n\to\infty]{} 0$, откъдето $x_n \xrightarrow[n\to\infty]{} x_0$.

\textbf{(в) $\implies$ (г):} Знаем, че съществува редица $\{x_n\}_{n=1}^{\infty} \subset D \setminus \{x_0\}$, $x_n \xrightarrow[n\to\infty]{} x_0$. Ще построим нейна подредица. Нека $n_1 := 1$. Тъй като $x_{n_1} \ne x_0$, имаме $|x_{n_1} - x_0| > 0$. Избираме $n_2 > n_1$ с $|x_{n_2} - x_0| < |x_{n_1} - x_0|$.

Нека сме избрали $n_1 < n_2 < \dots < n_k$ такива, че $|x_{n_i} - x_0| < |x_{n_{i-1}} - x_0|$ за всички $i \in \{2, \dots, k\}$. Отново $|x_{n_k} - x_0| > 0$ и от $x_n \xrightarrow[n\to\infty]{} x_0$ следва съществуването на $\bar{n} \in \mathbb{N}$ такова, че за всяко $n \ge \bar{n}$ е в сила $|x_n - x_0| < |x_{n_k} - x_0|$. Избираме $n_{k+1}$ такова, че $n_{k+1} \ge \bar{n}$ и $n_{k+1} > n_k$. Тогава $|x_{n_{k+1}} - x_0| < |x_{n_k} - x_0|$. По този начин подредицата $\{x_{n_k}\}_{k=1}^{\infty}$ удовлетворява условие (г).

\textbf{(г) $\implies$ (а):} Имаме редица $\{x_n\}_{n=1}^{\infty} \subset D \setminus \{x_0\}$ с $x_n \to x_0$ и $|x_{n+1} - x_0| < |x_n - x_0|$. Последните неравенства гарантират, че всички членове са попарно различни ($x_n \ne x_m$ при $n \ne m$). Тогава за произволно $\varepsilon > 0$ имаме:
\[
(x_0 - \varepsilon, x_0 + \varepsilon) \cap D \supset \{x_n : |x_n - x_0| < \varepsilon\} \supset \{x_n : n \ge n_0\},
\]
което е безкрайно множество.
\end{proof}

\section{Граница на функция}
\begin{definition}{Граница на функция (във формата на Коши)}{limit-cauchy}
Нека $f \colon D \to \mathbb{R}$, $D \subset \mathbb{R}$ и нека $x_0 \in \mathbb{R}$ е точка на сгъстяване на $D$. Казваме, че функцията $f$ има \textbf{граница} $L \in \mathbb{R}$, когато аргументът клони към $x_0$ (и пишем $\lim_{x\to x_0} f(x) = L$), ако за всяка околност $U$ на $L$ съществува околност $V$ на $x_0$ такава, че за всяко $x \in V \cap D$, $x \ne x_0$, е в сила $f(x) \in U$.

Еквивалентно на езика на $\varepsilon$--$\delta$:
\[
\lim_{x\to x_0} f(x) = L \iff \forall \varepsilon > 0 \; \exists \delta > 0 \; \forall x \in D \setminus \{x_0\} \colon (|x - x_0| < \delta \implies |f(x) - L| < \varepsilon).
\]
\end{definition}

\paragraph{Пример 3.6.}
Ще покажем, че $\lim_{x\to 0} \frac{\sin x}{x} = 1$. Дефиниционната област на $f(x) = \frac{\sin x}{x}$ е $(-\infty, 0) \cup (0, +\infty)$ и $0$ е нейна точка на сгъстяване.

\drawing{../figures/trig-inequality.pdf}{Тригонометрична окръжност и сравнение на лицата на триъгълник OAB, кръгов сектор OAB и триъгълник OAC}

Основните неравенства ще получим от геометрични съображения. Нека $x \in \left(0, \frac{\pi}{2}\right)$ е произволно. Означаваме с $O(0, 0)$ началото на координатната система, с $A(1, 0)$ пресечната точка на единичната окръжност с положителната част на абцисата, с $B(\cos x, \sin x)$ точка върху единичната окръжност, отговаряща на централен ъгъл $x$, и с $C(1, \operatorname{tg} x)$ пресечната точка на лъча $\vec{OB}$ с правата през $A$, перпендикулярна на абцисата.

Тогава лицето на $\triangle OAB$ (равно на $\frac{1}{2} \cdot 1 \cdot \sin x$) не надминава лицето на сектора $OAB$ от единичния кръг, което от своя страна не надминава лицето на $\triangle OAC$ (равно на $\frac{1}{2} \cdot 1 \cdot \operatorname{tg} x$). Тъй като лицето на сектора с централен ъгъл $x$ е $\frac{1}{2}x$, получаваме:
\[
\frac{1}{2} \sin x \le \frac{1}{2} x \le \frac{1}{2} \operatorname{tg} x \implies \sin x \le x \le \operatorname{tg} x = \frac{\sin x}{\cos x}.
\]
Тъй като всички количества са положителни за $x \in \left(0, \frac{\pi}{2}\right)$, следва:
\[
\cos x \le \frac{\sin x}{x} \le 1 \implies 1 - 2\sin^2 \frac{x}{2} \le \frac{\sin x}{x} \le 1
\]
\[
\implies 0 \le 1 - \frac{\sin x}{x} \le 2\sin^2 \frac{x}{2} \le 2\left(\frac{x}{2}\right)^2 = \frac{1}{2} x^2 \quad \text{за всяко } x \in \left(0, \frac{\pi}{2}\right).
\]
Поради четността на $\frac{\sin x}{x}$ и $x^2$, това неравенство е в сила за всяко $x \in \left(-\frac{\pi}{2}, \frac{\pi}{2}\right) \setminus \{0\}$:
\[
0 \le 1 - \frac{\sin x}{x} \le \frac{1}{2} x^2.
\]
Нека сега $\varepsilon > 0$ е произволно. Полагаме $\delta := \min\left\{\frac{\pi}{2}, \sqrt{2\varepsilon}\right\}$. Тогава за всяко $x \in (-\delta, \delta) \setminus \{0\}$ имаме:
\[
\left| \frac{\sin x}{x} - 1 \right| < \varepsilon \implies \lim_{x\to 0} \frac{\sin x}{x} = 1.
\]

\begin{definition}{Граница на функция (във формата на Хайне)}{limit-heine}
Нека $f \colon D \to \mathbb{R}$, $D \subset \mathbb{R}$ и нека $x_0 \in \mathbb{R}$ е точка на сгъстяване на $D$. Казваме, че функцията $f$ има \textbf{граница} $L \in \mathbb{R}$, когато аргументът клони към $x_0$, ако за всяка редица $\{x_n\}_{n=1}^{\infty} \subset D \setminus \{x_0\}$ от стойности на аргумента, клоняща към $x_0$, съответната редица от функционални стойности $\{f(x_n)\}_{n=1}^{\infty}$ клони към $L$:
\[
\left( \forall \{x_n\}_{n=1}^{\infty} \subset D \setminus \{x_0\}, \; x_n \xrightarrow[n\to\infty]{} x_0 \colon f(x_n) \xrightarrow[n\to\infty]{} L \right).
\]
\end{definition}

\begin{problembox}{Упражнение 3.8 (Задачи за обмисляне вкъщи)}
Нека $\{a_n\}_{n=1}^{\infty}$ е редица от реални числа. Докажете, че $a_n \xrightarrow[n\to\infty]{} \infty$ тогава и само тогава, когато $\{a_n\}_{n=1}^{\infty}$ няма точка на сгъстяване.
\end{problembox}